\documentclass[10pt,a4paper]{amsart}
\usepackage[margin=19mm]{geometry}
\usepackage{amsmath,amssymb,mathtools}
\usepackage[scr=rsfs,cal=euler]{mathalfa}
\usepackage{xfrac}
\usepackage[unicode,hidelinks,bookmarks=false]{hyperref}
\newcommand{\ie}{i.e.\ }
\newcommand{\cf}{cf.\ }
\newcommand{\resp}{respectively\ }
\newcommand{\Subsection}{\subsection}
\providecommand{\nobreakdash}{\nobreak\hbox{-}}
\setlength{\emergencystretch}{2em}
\multlinegap0pt
\usepackage{amssymb}
\usepackage{enumerate}
\newtheorem{lemma}{Lemma}
\title{Maximally almost periodic subgroups\\ of elementary Abelian topological groups$^\star$}
\author{Ol'ga Sipacheva}
\date{Source: arXiv:2605.17495v2 (CC BY 4.0)}
\begin{document}
\maketitle

\noindent\emph{Source and licence.} Maximally almost periodic subgroups\\ of elementary Abelian topological groups$^\star$. Version arXiv:2605.17495v2 (CC BY 4.0), distributed by arXiv under the Creative Commons Attribution 4.0 International licence, http://creativecommons.org/licenses/by/4.0/, as recorded in the arXiv licence metadata for this exact version and shown on the abstract page https://arxiv.org/abs/2605.17495v2. Full licence notice: \texttt{licenses/arxiv-2605.17495-CC-BY-4.0.txt}.\par\smallskip

\noindent\textit{Editorial scope.} Counted unit: the article's lemma l1 (source0.tex lines 228--237). The article's complete original proof of it is delivered with it (source0.tex lines 239--262, one \texttt{proof} environment of 24 lines). No mathematics is written, repaired or re-derived: the statement and the proof are the article's own lines, in the article's own notation, and no retained line is substituted or re-flowed.  No further source-local context is needed: every cross-reference of every retained line resolves inside the two delivered spans.  Nothing was omitted from any retained span, so the package carries no omission record.  No retained line contains a citation, so the wrapper carries no bibliography and no bibliography entry.  No retained line contains a figure environment, an image-inclusion command or drawing code, so no image is delivered and the package declares no asset dependency.  Editorial formatting only, in addition: an AMS \texttt{amsart} preamble that restates the article's own declarations and macros used by the retained lines, namely the environments \texttt{lemma} on separate counters, together with the standard packages those lines need.  The retained environments print in this document as Lemma 1 in order of appearance, whereas the article prints the same items with the numbering of its own full document.  The complete native source of the article as distributed in the arXiv e-print for this exact version is bundled unchanged as \texttt{sources/200799/source0.tex}.
\begin{lemma}
\label{l1}
Let $w=\lambda_1 e'_{i_1} + \lambda_2 e'_{i_2} +\dots + \lambda_n e'_{i_n}$\textup, 
where $n, i_1, \dots, i_n\in \mathbb N$\textup, $i_1< i_2< \dots < i_n$\textup, and $\lambda_1, \dots, 
\lambda_n\in \mathbb F_p\setminus \{0\}$\textup. Then 
$$ 
\|\mu e'_{i_{n-k}}\|\le p^k\cdot 2^k\cdot \|w\|\qquad \text{for each}\quad k=0, \dots, n-1\quad \text{and 
any}\quad \mu\in \mathbb F_p. 
$$ 
\end{lemma}

\begin{proof}
We argue by induction on $k$. Note that, for any $\mu \in \mathbb F_p$ and $x\in G$, we have $\|\mu x\|\le p\cdot 
\|x\|$ by virtue of the triangle inequality for norms. Thus, the required inequality for $k=0$ follows from property (iii) of $E'$. Suppose that 
$0<l<n$ and the inequality holds for any $k<l$ and $\mu\in \mathbb F_p$. Applying the triangle inequality for 
norms, we obtain 
\begin{multline*}
\|(p-\lambda_{n-l+1})e'_{i_{n-l+1}}+(p-\lambda_{n-l+2})e'_{i_{n-l+2}}+\dots +(p-\lambda_{n})e'_{i_{n}}\| \\
\le p^{l-1}\cdot 2^{l-1}\cdot\|w\|+ p^{l-2}\cdot 2^{l-2}\cdot\|w\|+\dots+p^0\cdot 2^0\cdot\|w\|\\
\le p^{l-1}\cdot (2^{l-1}+2^{l-2}+\dots+1)\cdot\|w\|= p^{l-1}\cdot (2^l-1)\cdot \|w\|. 
\end{multline*}
Noting that 
\begin{multline*}
\lambda_1 e'_{i_{1}}+\lambda_2 e'_{i_{2}}+\dots +\lambda_l e'_{i_{l}}  \\ = w + 
(p-\lambda_{n-l+1})e'_{i_{n-l+1}}+(p-\lambda_{n-l+2})e'_{i_{n-l+2}}+\dots +(p-\lambda_{n})e'_{i_{n}}
\end{multline*}
and 
applying the triangle inequality once more, we arrive at 
$$
\|\lambda_1 e'_{i_{1}}+\lambda_2 e'_{i_{2}}+\dots +\lambda_l e'_{i_{l}}\|\le 
\|w\|+ p^{l-1}\cdot (2^l-1)\cdot \|w\| \le p^{l-1}\cdot 2^l\cdot\|w\|. 
$$
By property (iii) we have $\|e'_{i_{l}}\|\le p^{l-1}\cdot 2^l\cdot\|w\|$ and $\|\mu e'_{i_{l}}\|\le p^l \cdot 
2^l\cdot\|w\|$. 
\end{proof}

\end{document}
